\chapter{Conclusion}
\label{ch:conclusion}

\begin{aside}
You have arrived at the end. Eight hundred pages on a library
that prints text. Worth it? Up to you.
\end{aside}

\section{What you've learned}

The terminal as a machine. \cpp{}26 features. \maya{}'s DSL,
elements, layout, rendering pipeline, signals, runtime,
widgets. Patterns. Anti-patterns. Build systems. Distribution.
Future directions.

That's a stack. And a TUI is not a tiny thing once you look
inside.

\section{What's next}

Build a thing. Read the source. Contribute. The maintainers
are friendly; the codebase is small; the niche is alive.

\section{Thank you}

For your time. For your attention. For your willingness to
believe that text-on-screen is interesting enough to read
about.

\section{A last note}

The terminal is older than most readers of this book. It has
outlasted many GUI fashions. It will outlast many more. The
character grid is a stable medium for stable ideas.

\maya{} is one library in one moment of that long arc. Make
something with it. Or with whatever comes next.

\medskip

Goodbye.
